% Lösungen Einheit 2 von 4 – Sinus- und Kosinusfunktion und ihre Merkmale (zusammenbau v0.1)
\einheitenkopf{Einheit 2 von 4 · Sinus- und Kosinusfunktion und ihre Merkmale}

\erg{16}{a) von A bis C \quad b) 0; 0,71; 1; 0,71; 0; −0,71; −1; −0,71; 0 \quad c) runde Welle durch (0|0), (90|1), (180|0), (270|−1), (360|0), waagerecht im Hoch- und Tiefpunkt \quad d) 40° je cm; beschriften etwa 0°, 80°, 160°, 240°, 320° – in Grad, nicht in cm \quad e) sin 135° ≈ 0,7 \quad f) x ≈ 45° und x ≈ 135° \quad g) ja: sin 50° ≈ 0,77 \quad h) Grenzen −1 und 1: alle Werte von −1 bis 1 \quad i) 0; 180; 360; 540; 720 (in Grad) – alle 180° eine Nullstelle \quad j) H(90|1); T(270|−1) \quad k) 360° – dann läuft der Punkt ein zweites Mal um den Kreis \quad l) fällt – der Abschnitt liegt zwischen Hochpunkt und Tiefpunkt \quad m) punktsymmetrisch zum Ursprung: sin(−40°) = −sin 40° ≈ −0,64 \quad n) 90° \quad o) y = cos(x) \quad p) Werte 0; −0,71; −1; −0,71; 0; 0,71; 1; 0,71; 0; Wertebereich von −1 bis 1; Nullstellen bei −180°, 0°, 180°; H(90|1), T(−90|−1); kleinste Periode 360°; punktsymmetrisch zum Ursprung}

16c)

\begin{ksys}[xmin=0,xmax=360,ymin=-1.2,ymax=1.2,xstep=45,ystep=0.2,klein]\funktion{sin(\x)}{f}\end{ksys}


16p)

\begin{ksys}[xmin=-180,xmax=180,ymin=-1.2,ymax=1.2,xstep=45,ystep=0.2,klein]\funktion{sin(\x)}{f}\end{ksys}

\erg{17}{a) Fehler: Die Welle hat oben keine Spitze. Richtig: rund gebogen und waagerecht im Hochpunkt H(90|1) \quad b) Nach einer vollen Drehung steht der Punkt auf dem Einheitskreis wieder am selben Ort; ab dort läuft er dieselben Höhen noch einmal durch. \quad c) g – bei 0° ist der Wert 1}
